\subsection{由指数得到的秩界}\label{sec:finite-ranks}

我们将从固定规模的算法过渡到张量幂上的算法，并结合边界秩的闭包论证~\cite{bini,bcs,spectrum}。

\begin{proof}[推论~\ref{cor:tight} 的证明]
固定 \(n\ge2\)。对于任意域 \(K\)，若 \(T_{n^j}\) 在 \(K\) 上有一个秩为 \(r\) 的分解，由引理~\ref{lem:specialize} 中的递归与填充，有 \(\omega_K\le\log_{n^j} r\)。因此由定理~\ref{thm:square}，对于每个 \(j\)，均有 \(\operatorname{rank}_K(T_{n^j})\ge n^{9j/4}\)。取 \(j\) 次方根并令 \(j\) 趋于无穷，便得渐近秩的下界。令 \(K=\overline F\) 且 \(j=1\)，由引理~\ref{lem:bilinearize} 即得所述非线性门数的下界 \(\lceil n^{9/4}/2\rceil\)。

为证明几何边界秩的下界，在 \(\overline F\) 上包含 \(T_n\) 的固定张量空间 \(V\) 中工作。设 \(T_n\) 属于所有秩至多为 \(r\) 的张量的 Zariski 闭包。对于每个 \(j\)，这些张量的 \(j\) 次幂所张成的空间包含 \(T_n^{\otimes j}\)，因为张量幂属于该线性空间是一个闭条件。关于 \(\dim V\) 个变量的 \(j\) 次齐次多项式的系数个数给出了该空间维数的上界。因此 \(T_n^{\otimes j}\) 可以表示为 \(O_n(j^{\dim V-1})\) 个这种幂的线性组合，每个幂的秩至多为 \(r^j\)。与阶数 \(n^j\) 处的下界比较，再取 \(j\) 次方根并令 \(j\) 趋于无穷，便得 \(r\ge n^{9/4}\)。

为证明渐近秩的上界，固定 \(\varepsilon>0\)。在阶数 \(n^j\) 处，由已知的电路上界及引理~\ref{lem:bilinearize}，可得 \(\overline F\) 上的一个分解，其秩为 \(r_j=O_{F,\varepsilon}(n^{j(9/4+\varepsilon)})\)。这些系数属于某个次数为 \(d_j\) 的有限扩张 \(E_j/F\)。对每个 \(t\)，将此分解取 \(t\) 次张量幂。在 \(E_j\) 的一组包含 \(1\) 的 \(F\)-基下展开每个简单张量的三个向量，再将展开后出现在各项系数中的乘积投影到基向量 \(1\) 所对应的坐标上。所表示的张量定义在 \(F\) 上，故得
\[
        \operatorname{rank}_F(T_n^{\otimes jt})\le d_j^3 r_j^t .
\]
由次乘性，定义渐近秩的极限存在。先取 \(jt\) 次方根，固定 \(j\) 令 \(t\to\infty\)，再令 \(j\to\infty\) 及 \(\varepsilon\downarrow0\)。由此在原域上得到上界 \(n^{9/4}\)。
\end{proof}

在有限阶数下，推论~\ref{cor:tight} 还与 Sun 的预印本第一版~\cite{Sun} 所提出的关于二的幂次阶数的边界秩估计相冲突。当 $n=2^8$ 时，该文提出的上界为 $2^{17}$，而我们的下界为 $2^{18}$。

方阵乘法张量还判定了关于紧三阶张量的一项预测。若每个局部基都可赋予两两不同的整数权重，且在支撑中的每个三元组上相应权重之和为零，就称该支撑是紧的。在 \(T_n\) 中，取 \(B>n\)，分别给 \(e_{ij},e_{jk},e_{ki}\) 赋予权重 \(i-Bj\)、\(Bj-B^2k\) 和 \(B^2k-i\)。在各自的因子中，每个赋权映射都是单射，且三项权重之和为零。沿任意单个因子展平所得的秩均为 \(n^2\)。在这种情形下，Strassen 的紧张量猜想预测渐近秩为 \(n^2\)~\cite[猜想~5.3]{Strassen1994}。推论~\ref{cor:tight} 给出的则是 \(n^{9/4}\)。

\begin{proof}[推论~\ref{cor:rectangular} 的证明]
沿三个因子分别展平所得的秩为 \(ab,bc,ca\)。由于矩阵秩的上界条件是闭条件，这三个数均为几何边界秩的下界。由于展平秩在张量积下可乘，它们也是渐近秩的下界。在
\[
            T_{a,b,c}\otimes T_{b,c,a}\otimes T_{c,a,b}
                                \simeq T_{abc}
\]
中将对应因子归组，三个指标中的每一个都变为一个有 \(abc\) 种可能取值的元组。这两种秩在因子的循环重标号下不变，且对张量积具有次乘性。因此，由推论~\ref{cor:tight}，二者在 \(T_{a,b,c}\) 处的值各自取立方，都至少为 \((abc)^{9/4}\)。\(abc=1\) 的情形显然成立。

当 \(a=c=n\) 且 \(b=\lfloor n^\beta\rfloor\) 时，几何边界秩的下界为 \(\Omega(n^{3(2+\beta)/4})\)。引理~\ref{lem:bilinearize} 中的双线性化同样可逐项用于矩形矩阵乘法。因此，若指数为二，就必须有 \(3(2+\beta)/4\le2\)。由此得 \(\alpha_F\le2/3\)，而 \(\beta=1/2\) 时的乘法上界给出 \(\alpha_F\ge1/2\)~\cite{abvz}。
\end{proof}
